\chapter{Conclusion}
\label{chap:conclusion}

% ============================================================
\section{The Research Problem Revisited}
\label{sec:ch5_problem}
% ============================================================

This dissertation began from a practical observation: corporate social responsibility is easy to declare and difficult to implement. A firm can announce a take-back commitment, a recycling target, or a responsible-sourcing policy at little immediate cost. Delivering on that announcement is a different matter. It requires someone to invest in capability, someone to operate a physical return channel, someone to qualify treatment partners, and some rule by which the resulting cost is divided among firms that remain legally and financially independent.

The difficulty is that these roles rarely sit with the same firm. The party that benefits most from a responsible brand position is often not the party that bears the operational burden of making the position credible. A Manufacturer may own the brand, the product design, and the regulatory obligation, while retailers operate the consumer-facing return points and absorb the handling cost. When responsibility and benefit are separated in this way, CSR stops being only an ethical question and becomes a governance and coordination question: who decides, who pays, who gains, and what holds the arrangement together.

Chapter~\ref{chap:introduction} framed this as the dissertation's central problem and posed one overarching question with two study-specific questions. Chapter~\ref{chap:literature_review} showed that the surrounding literatures each address part of the problem but leave the connection between them underdeveloped. Chapters~\ref{chap:model1} and~\ref{chap:model2} then developed two complementary models: one for the bilateral case, where the question is who controls and finances CSR maturity, and one for the network case, where the question is how several firms build shared stewardship infrastructure and divide its cost. This chapter draws the two together.

% ============================================================
\section{Findings from Study 1: Bilateral CSR Governance}
\label{sec:ch5_study1}
% ============================================================

Chapter~\ref{chap:model1} modeled a single Manufacturer selling to a single Retailer, with the Manufacturer setting a wholesale price, the Retailer setting a markup and a level of CSR maturity, and consumer demand rising with CSR maturity but falling with the final retail price. Four governance regimes were compared: Vertical Nash, Manufacturer Stackelberg, Retailer Stackelberg, and an integrated benchmark.

Three findings carry beyond the specific functional forms.

First, \emph{governance structure changes CSR outcomes even when technology and preferences do not}. The same demand parameters and the same cost of CSR investment produce different equilibrium CSR maturity, different prices, and different profit splits depending only on who moves first. Integration consistently delivers the highest total profit, demand, and CSR maturity because it removes double marginalization; the decentralized regimes fall short by margins that depend on the channel's power structure. A firm that wants to explain weak CSR performance therefore cannot look only at the cost of the CSR investment. The allocation of pricing authority is itself a determinant of the outcome.

Second, \emph{the best decentralized arrangement depends on how strongly consumers respond to CSR}. The cross-regime comparison identified two parameter regions. Under moderate CSR pull, total channel profit ranks Vertical Nash above Manufacturer Stackelberg, making simultaneous play the second-best decentralized structure. Under high CSR pull, the ranking reverses and Manufacturer leadership becomes second-best, because a leading Manufacturer supports the higher CSR investment that strongly responsive consumers reward. Retailer Stackelberg is the weakest structure in both regions. There is consequently no single governance recommendation that holds across markets; the recommendation depends on a measurable feature of demand.

Third, \emph{coordination is itself an equilibrium outcome and can fail}. The Stage-0 meta-game treated the choice between coordinating and competing for leadership as a strategic decision made before any operational decision. Under moderate CSR pull the resulting game is a stag hunt with two equilibria, so the efficient coordinated outcome is available but not guaranteed: it requires mutual confidence, and firms can rationally settle on the inferior outcome. Only under high CSR pull does coordination become the unique equilibrium. This is the chapter's most policy-relevant result. Demonstrating that integration is jointly efficient does not establish that independent firms will reach it, and interventions that raise consumer CSR responsiveness can change the equilibrium structure rather than merely shifting payoffs within it.

% ============================================================
\section{Findings from Study 2: Collective Product-Stewardship Governance}
\label{sec:ch5_study2}
% ============================================================

Chapter~\ref{chap:model2} changed the unit of analysis. One Manufacturer and several independent retailers must jointly build shared infrastructure, qualify treatment partners, route returned products, and divide the resulting physical cost. The chapter used a two-phase structure: a coalition-dependent mixed-integer program that determines what each possible coalition could achieve, and a cooperative cost game built on those coalition costs.

Four findings stand out.

First, \emph{a stewardship promise must be converted into an operating requirement before it can be costed}. The formal-service target states how much returned volume the coalition commits to handle through the qualified network. Because it is fixed in advance rather than optimized, the model cannot generate savings by quietly reducing service. This reframes the question from how much a firm claims to care about responsibility to what service it has committed to and whether the network can deliver it.

Second, \emph{qualified partner access has option value rather than guaranteed savings}. Adding an eligible processor cannot increase a coalition's optimized cost, but the option matters only when the network economics select it. Certification, due diligence, and contracting create a feasible choice; they do not by themselves reduce cost.

Third, \emph{shared infrastructure makes stewardship economics nonlinear}. Fixed facilities create scale economies, but capacity thresholds can force an additional center or processor. A member's marginal contribution therefore depends on the order in which coalitions form, not only on the volume it brings. This is the operational justification for allocating cost by the Shapley value rather than by volume share.

Fourth, and most importantly, \emph{cost sharing and coalition stability are sequential and distinct questions}. In the reproducible case the shared network reduced net stewardship cost from \ChFourStandaloneTotal\ to \ChFourGrandCost\ thousand U.S. dollars, a saving of \ChFourSavingsPercent. Every player individually preferred participating to operating alone. Yet the Shapley allocation failed the full core test, with a minimum slack of \ChFourMinSlack\ thousand U.S. dollars: at least one subgroup of members was assigned more than it would pay by leaving and operating together. The least-core value of \ChFourLeastCoreEpsilon\ thousand U.S. dollars quantifies the smallest blocking incentive that no budget-balanced allocation can remove.

A structural diagnosis then explains the failure. If the network's facility decisions were divisible, the core would always be nonempty, so the instability comes entirely from facilities that must be opened whole. A small balanced collection of subgroups, most of them led by the Manufacturer, proves that the core is empty; it can be named and checked by hand. These subgroups can run on the high-standard processor alone, while the full coalition, with more volume than that processor can carry, must also activate a second one. The instability therefore appears only in a predictable regime. In the case, lowering the formal-service target to the point where one processor suffices restores stability. The instability also has a price: an external contribution of \ChFourCostOfStability\ thousand U.S. dollars, \ChFourCoSShareSavings\ of the savings cooperation creates, would remove every blocking incentive.

That contrast is the chapter's central result and it generalizes well beyond the case data. Budget balance is a property of the allocation formula and holds automatically. Individual participation compares each firm with its own outside option and is a weak test, since it examines only single-member deviations. Coalition stability examines every subgroup and is strictly stronger. An allocation can satisfy the first two and fail the third, and in this case it does.

The chapter also examined governance-motivated alternatives to pure marginal-contribution allocation, including a hybrid rule that blends the Shapley value with an explicit responsibility share, a budget-neutral responsibility transfer to the Manufacturer, and a least-core projection. These produced a result worth stating plainly: allocations that shift more cost onto the Manufacturer on responsibility grounds did not improve stability and in several cases made the worst subgroup's blocking incentive larger. Aligning an allocation more closely with a responsibility principle is therefore not the same as making it more acceptable to the parties who must sustain it. Whether a rule improves stability is an empirical question that must be recomputed, not an inference from the rule's normative appeal.

% ============================================================
\section{Answers to the Research Questions}
\label{sec:ch5_answers}
% ============================================================

The dissertation can now answer the questions posed in Chapter~\ref{chap:introduction}.

\paragraph{Overarching question: how can CSR commitments be operationalized in supply chains so that independent firms have incentives to invest, coordinate, and share responsibility?}
The two studies together indicate that this requires solving two different problems that are frequently treated as one. The first is a control-and-return problem: some party must have both the authority to act and a sufficient share of the resulting benefit to justify the investment. The second is a distribution problem: once a shared capability exists, its cost must be divided in a way that no participant, and no group of participants, can improve upon by withdrawing. Solving either problem alone is insufficient. A well-governed bilateral channel that cannot scale to shared infrastructure leaves system-level responsibility unaddressed, while an efficient shared network with an unstable cost-sharing rule invites the defection that dissolves it.

\paragraph{Study 1 question: in a bilateral channel, how do vertical power, pricing timing, and a stage-0 regime-choice game affect CSR-driven maturity, demand, and profit allocation?}
Vertical power and timing affect all three, and they do so systematically rather than marginally. Integration dominates on CSR maturity, demand, and total profit; Retailer leadership is consistently weakest because it shifts margin downstream while suppressing the CSR investment that generates demand. Between the intermediate structures, the ranking switches with the strength of consumer CSR response, so that simultaneous play is preferable when responsiveness is moderate and Manufacturer leadership is preferable when it is strong. The Stage-0 analysis adds that the efficient regime is not always selected: under moderate CSR pull, coordination and leadership competition are both equilibria, and the outcome turns on expectations rather than on payoffs alone.

\paragraph{Study 2 question: in a multi-retailer product-stewardship network, how should shared infrastructure be designed, how should cost be allocated so that every firm prefers participating, and does that allocation survive the stronger test that no subgroup prefers to leave?}
The design question is answered by the coalition-dependent network model, which shows that the cheapest configuration depends on which members are present, what service level they have committed to, and which treatment partners they may use. The allocation question is answered affirmatively for individual participation: Shapley allocation fully funded the network and left every firm better off than acting alone. The stability question is answered negatively for the case examined. The allocation was blocked by subgroups, and the remedies considered reduced but did not eliminate the blocking incentive. The structural diagnosis turns this negative answer into an explanation. The failure is not a defect of the Shapley rule, since no allocation removes it; it arises because indivisible facilities let certain Manufacturer-led subgroups avoid a cost the full coalition cannot. It occurs only when the service commitment exceeds what a single qualified processor can carry, and differences between retailers do not cause it. The honest answer is therefore that the first two parts are solved and the third is not solved automatically. The model's contribution is to make the shortfall visible, explain its cause, and price its removal, rather than to conceal it.

% ============================================================
\section{Cross-Study Synthesis}
\label{sec:ch5_synthesis}
% ============================================================

The two studies are usually read as a progression from a small model to a larger one. The more useful reading is that they identify two independent failure modes, either of which is sufficient to prevent a CSR commitment from being implemented.

Chapter~\ref{chap:model1} describes an \emph{incentive-to-invest} failure. Even with only two firms and complete information, the party best positioned to build CSR capability may not be the party that captures its returns, and the firms may fail to coordinate on the efficient governance regime even when both would prefer it. This failure occurs before any physical infrastructure exists.

Chapter~\ref{chap:model2} describes a \emph{willingness-to-remain} failure. Even when shared infrastructure is efficient, fully funded, and individually attractive to every participant, a subgroup may still find withdrawal profitable. This failure occurs after the infrastructure exists. It is caused by the structure of the network's costs rather than by the choice of allocation rule: when facilities are indivisible, no division of the cost can satisfy every subgroup at once.

The two failures are structurally different and call for different instruments. The first is addressed by changing who leads, who invests, or how strongly consumers reward responsible behavior; contractual and demand-side interventions apply. The second cannot be solved by reallocating cost among the members alone. It is addressed by altering the outside options that make defection attractive, by investing in shared governance that any breakaway group would have to duplicate, by redesigning capacity, or by an external contribution. The case offers two illustrations. Removing access to the Manufacturer-enabled high-standard processor removed the blocking incentive entirely, at the price of a slightly more expensive network. The reason is that the Manufacturer-led subgroups could no longer run on that processor alone while the full coalition needed a second one. Lowering the formal-service target far enough had the same effect. This exposes a tension that matters for CSR policy: a stronger stewardship commitment can be exactly what makes a collective unstable. Efficiency, ambition, and stability are not always complementary, and a governance design may reasonably accept a marginally costlier or less ambitious arrangement in exchange for one that holds together---or fund the modest contribution that keeps the more ambitious one intact.

The unifying claim of the dissertation is therefore this: credible CSR in a supply chain requires both a governance regime that gives some party the incentive to invest, and a cost-sharing rule that independent parties will not abandon. Neither condition implies the other, and neither alone is sufficient. Frameworks that treat CSR implementation only as a matter of moral commitment, only as a matter of operational efficiency, or only as a matter of fair-sounding cost allocation each address a fragment of the problem.

% ============================================================
\section{Contributions}
\label{sec:ch5_contributions}
% ============================================================

The dissertation makes four contributions, stated at the scope the analysis supports.

First, it operationalizes CSR as a measurable capability rather than a preference. In Chapter~\ref{chap:model1} CSR appears as a costly maturity investment that shifts demand; in Chapter~\ref{chap:model2} it appears as a formal-service commitment, qualified-partner access, and governance expenditure. In both cases CSR enters the model through decisions a firm actually makes.

Second, it treats the governance regime as an object of analysis rather than an assumption. The Stage-0 meta-game makes the choice between coordination and leadership competition an equilibrium outcome, which allows the dissertation to ask why coordination fails rather than to assume that it succeeds.

Third, it provides a transparent operational bridge between network design and cooperative allocation. Prior work has combined optimization with cost allocation; the specific contribution here is a formulation in which a service commitment, manufacturer-enabled partner access, and facility decisions jointly determine each coalition's realizable cost, together with a fully enumerated and reproducible stability diagnosis. That diagnosis goes beyond reporting a failure: it proves that instability can arise only from indivisible facilities, names the responsible subgroups through a balanced-collection certificate, predicts how stability responds to the network inputs, and bounds the external contribution needed to restore it.

Fourth, and most transferable, it separates four questions that are often conflated: whether an allocation funds the network, whether each firm prefers to join, whether every subgroup prefers to stay, and whether the arrangement is normatively fair. The first is guaranteed by construction, the second and third must be computed and can diverge, and the fourth is not settled by any of the preceding three. Keeping these apart is what allows the dissertation to report a system that saves \ChFourSavingsPercent\ and is individually attractive while still being coalitionally unstable, instead of presenting a single undifferentiated claim of fairness.

% ============================================================
\section{Limitations and Directions for Future Research}
\label{sec:ch5_limitations}
% ============================================================

The models are deliberately stylized, and their boundaries indicate where further work is most valuable.

\paragraph{Model structure.} Chapter~\ref{chap:model1} uses a single product, deterministic demand, complete information, and a one-shot Stage-0 interaction. Chapter~\ref{chap:model2} is deterministic, linear in variable costs, single-period, and treats coalition formation as one-shot. Real stewardship systems face uncertain return volumes, changing battery chemistries, volatile commodity prices, asymmetric information, and repeated interaction. Stochastic and robust formulations of the network phase, and repeated-game formulations of the allocation phase, are natural extensions; repetition in particular may sustain cooperation that a one-shot core test rejects.

\paragraph{Empirical calibration.} This is the most immediate gap. The Chapter~\ref{chap:model2} case is reproducible and fully auditable, but its inputs are stylized rather than estimated. Regional return volumes, facility and transport costs, processor contract terms, governance expenditure, and coalition outside options all require empirical grounding. Some of these are publicly documented; facility economics and processor contract terms are commercially sensitive and may require interviews or negotiated data access. Until this work is done, the numerical results demonstrate mechanism rather than magnitude.

\paragraph{Measurement of the model's inputs.} The formal-service target, partner eligibility, and governance cost are observable in principle but were set illustratively here. Service commitments could be estimated from published program rules, eligibility verified from contracts and certifications, and governance spending measured directly. Interviews with stewardship managers would also clarify how these terms are negotiated before network design begins, which the present model treats as given.

\paragraph{Outside options and the Manufacturer's role.} Coalition outside options are modeling assumptions with substantial influence on the stability results. The Manufacturer's stand-alone cost is interpreted as platform readiness rather than treatment of physical volume, and this interpretation drives its comparatively low Shapley allocation. Alternative interpretations should be tested. Similarly, recovery performance is currently reported rather than constrained; imposing a binding recovery target would require re-solving every coalition, since both feasibility and outside options would change.

\paragraph{Beyond economic acceptance.} The dissertation's stability tests are economic. They compare modeled payments with modeled outside options and say nothing about transparency, voice, bargaining power, legal responsibility, or procedural consistency. The finding that responsibility-weighted allocations worsened measured stability raises a question this dissertation cannot settle: whether such allocations nonetheless improve perceived legitimacy, and whether legitimacy sustains cooperation that pure economic incentives would not. That question is behavioral and institutional, and it is the natural continuation of this work.

\paragraph{Computational scope.} Exact enumeration of all coalitions is tractable for the six players studied here but grows exponentially. Larger stakeholder sets will require sampling, decomposition, or approximation methods, and the accuracy of those approximations for stability testing is itself an open question.

These limitations share a structure that reflects the design of the framework. New operational detail can be added to the network phase without disturbing the allocation phase, and new allocation rules can be tested without rebuilding the network model. The two-phase separation is intended to remain stable as the empirical content improves.
